\chapter{Reduction of Endomorphisms}

\section{Principal Ideal Domain}

\begin{definitionenv}
    Let $R$ be a commutative unitary ring.
    By definition, an ideal of $R$ is a sub-$R$-module of $R$. 
    If an ideal $I$ of $R$ is generated by one element, we say that $I$ is a \textbf{principal ideal}.

    \quad If all ideals of $R$ are principal, we say that $R$ is a \textbf{principal ideal ring}.

    \quad Let $a,b \in R$, then $aR \subseteq bR$ if and only if there exists $c \in R$ such that $a = b c$.
    We denote this condition as $b \mid a$, and say that $b$ is a \textbf{divisor} of $a$ and $a$ is a \textbf{multiple} of $b$.
\end{definitionenv}

\begin{definitionenv}\label{prop:6.1.3}
    Let $R$ be a commutative unitary ring.
    If the neutral element for the addition and multiplication are different, and if
    $$\forall s,t \in R\setminus\{0\},\ st \in R\setminus \{0\},$$
    then we say that $R$ is an \textbf{integral domain}.
\end{definitionenv}

\begin{propositionenv}
    Let $R$ be an integral domain and $a,b \in R$.
    Then $aR = bR$ if and only if there exists $c \in R^\times$ such that $a = b c$.
\end{propositionenv}
\begin{proofenv}
    We assume that there exists $c \in R^\times$ such that $a = b c$.
    By definition, $aR \subseteq bR$.
    Since $c$ is invertible, also $bR \subseteq aR$.

    \quad Conversely, assume that $aR = bR$, then we have $x,y \in R$ such that $a = b x$, $b = a y$.
    If $a = 0$, then $b = 0$.
    In this case, $a = 1 \cdot b$.
    If $a \neq 0$, then $a = b x = a y x$.
    So $y x = 1$ and $x y \in R^\times$.
\end{proofenv}

\begin{definitionenv}
    Let $K$ be a commutative unitary ring.
    Recall that a non-zero polynomial $P \in K[T]$ can be written as
    $$a_0 + a_1 T +\cdots + a_n T^n,$$
    where $a_0,\cdots a_n \in K$ and $a_n \neq 0$.
    
    \quad The natural number $n$ is called the \textbf{degree} of $P$, and denote as $\deg (P) \in \NN$.
    The element $a_n$ is called the \textbf{leading coefficient} of $P$, and denote as $c(P) \in K$.

    \quad If $c(P) = 1$, then we say that $P$ is a \textbf{monic}.

    \quad By convention, the degree of $0$ is defined to be $-\infty$,
    and the zero polynomial is assumed to be monic.
\end{definitionenv}

\begin{propositionenv}
    Let $K$ be an integral domain.
    Then the ring $K[T]$ is also a integral domain.
    It is a principal ideal ring if $K$ is a field.
\end{propositionenv}
\begin{proofenv}
    Since $K$ is an integral domain, if $P,Q \in K[T]$ are non-zero in $K[T]$, then
    $$\deg (PQ) = \deg (P) + \deg (Q),\quad c(PQ) = c(P)c(Q).$$
    Hence $K[T]$ is also an integral domain.

    \quad We check that $K[T]$ is a principal ideal ring when $K$ is a field.
    We solve it using Euclidean division of polynomials.
    If $I$ is a zero ideal, then $(0) = I$.
    If $I$  is non-zero, then there exists non-zero $P \in I$ such that $\deg(P)$ is minimal inside $I$.
    For any $F \in K[T]$, there exists a unique pair $(Q,R) \in K[T]^2$ such that 
    $$F = PQ + R, \quad \deg(R) < \deg(P).$$
    Hence, since the degree of $P$ is minimal, $R = 0$, so $F = PQ$.
    So $I = (P)$.
\end{proofenv}

\begin{remark}
    Let $K$ be a field,
    and $I$ be an ideal of $K[T]$.
    Then there exists a unique monic $P \in I$ such that $I = (P)$.
\end{remark}

\begin{definitionenv}
    Let $A$ be an integral domain.
    We say that $a \in A\setminus A^\times$ is \textbf{irreducible} if $a$ can not be decomposed as a product of two elements in $A\setminus A^\times$.

    \quad In case when $A = K[T]$,
    an irreducible element is called an \textbf{irreducible polynomial}.
\end{definitionenv}

\begin{definitionenv}
    Let $A$ be a commutative unitary ring and $I$ be an ideal of $A$.
    We say that $I$ is \textbf{prime ideal} if the quotient ring $A/I$ is an integral domain.
    This is equivalent to the condition that $I \neq A$ and
    $$\forall a,b \in A \setminus I, ab \in A \setminus I.$$
    We say that $I$ is \textbf{maximal} if $A/I$ is a field.
\end{definitionenv}

\begin{propositionenv}\label{prop:6.1.9}
    Let $A$ be a commutative unitary ring.
    The maximal ideals of $A$ are precisely the maximal elements in the set of al ideals of $A$ ordered by inclusion $\subseteq$.
\end{propositionenv}
\begin{proofenv}
    Let $I$ be an ideal of $A$.
    Suppose that $A/I$ is a field.
    If $J$ is another ideal of $A$ such that $J \supsetneq I$,
    then there exists $a \in J \setminus I$ such that the class of $a$ in $A/I$ is non-zero.
    This implies that there exists $b \in A$ such that $ab - 1 \in I$.
    In particular,
    $$1 = (1 - ab) + ab \in J.$$
    Hence, $J = A$.

    \quad Conversely, suppose that $I \neq A$ is maximal amongst ideals of $I$ with respect to inclusion.
    If $a \in A \setminus I$ is such that the class of $a$ is not invertible in $A/I$, then the set
    $$J = \{ab + c \mid b \in A, c \in I\}$$
    is an ideal of $A$ containing $I$.
    $J \neq A$, since $1$ does not belongs to $J$.
    Contradiction with the maximality of $I$.
\end{proofenv}

\begin{propositionenv}
    Let $A$ be an integral domain, $a \in A$.
    If $(a)$ is a prime ideal,
    then $a$ is an irreducible element of $A$.
\end{propositionenv}
\begin{proofenv}
    Suppose that $a = bc$ with $b,c \in A\setminus A^\times$.
    Since $I = (a)$ is prime,
    then $b$ and $c$ can't be simultaneously belong to $A\setminus I$. 
    Suppose that $b \in I$,
    then there exists $d \in A$ such that
    $b = ad = bcd$.
    Hence $cd = 1$.
    This contradicts $c \in A \setminus A^\times$.
\end{proofenv}

\begin{propositionenv}
    Let $A$ be a principal ideal domain.
    If $a \in A$ is irreducible,
    then the ideal generated by $a$ is a maximal ideal.
\end{propositionenv}
\begin{proofenv}
    Let $I = (a)$.
    By proposition \ref{prop:6.1.9},
    it suffices to show that $I$ is not contained in any other ideal.

    \quad Suppose by contradiction that $I \subsetneq J \subsetneq A$.
    Since $A$ is a principal ideal domain,
    there exists $b \in A$ such that $J = (b)$.
    Since $I \subseteq J$,
    there exists $c \in A$ such that $a = bc$.
    Notice that $b \in A \setminus A^\times$, $c \in A \setminus A^\times$ since otherwise we would have $I = J$.
    Contradiction.
\end{proofenv}

\begin{definitionenv}
    Let $R$ be a commutative unitary ring.
    We say that $R$ is \textbf{noetherian} if for any increasing sequence
    $$I_0 \subseteq I_1 \subseteq I_2 \subseteq \cdots$$
    of ideals in $R$, there exists $n \in \NN$ such that
    $$I_n = I_{n+1} = I_{n+2}.$$
\end{definitionenv}

\begin{propositionenv}
    A commutative unitary ring $R$ is noetherian if and only if any ideal of $R$ is finitely generated.
    Any principal ideal domain is noetherian.
\end{propositionenv}
\begin{proofenv}
    Assume $R$ is noetherian.
    Let $I$ be an ideal of $R$.
    Assume that $I$ is not finitely generated.
    One can recursively construct a strictly increasing sequence
    $$I_0 \subseteq I_1 \subseteq \cdots I_n \subseteq I_{n+1} \subseteq \cdots$$
    using the infinite set of generations of $I$.
    Contradiction.
    Conversely assume that any $I$ in $R$ is finitely generated.
    Let 
    $$I_0 \subseteq I_1 \subseteq \cdots I_n \subseteq I_{n+1} \subseteq \cdots$$
    be an increasing sequence of ideals.
    Then $\displaystyle I \coloneq \bigcup_{n \in \NN} I_n$ is an ideal in $R$, so $I = (x_1,\cdots,x_p)$.
    Let $n \in \NN$ such that
    $$I_n \supseteq \{x_1,\cdots x_p\},$$
    hence $(x_1,\cdots,x_p) = I_n = I_{n+1} = \cdots$.
\end{proofenv}

\begin{corollaryenv}\label{cor:6.1.14}
    Let $R$ be a commutative unitary noetherian ring.
    Then any ideal $I \neq R$ is contained in some maximal ideal of $R$.
\end{corollaryenv}
\begin{proofenv}
    Suppose that $I$ is not contained in any maximal ideal, then recursively we obtain a strictly increasing sequence
    $$I \subseteq I_0 \subsetneq I_1 \subsetneq \cdots \subsetneq I_n \subsetneq I_{n+1} \subsetneq \cdots$$
    such that $I_n \neq R$ for any $n \in \NN$.
\end{proofenv}

\begin{remark}\label{6.1.15}
    The statement of corollary \ref{cor:6.1.14} is true for any commutative unitary ring under the axiom of choice.
    Let $R$ be an commutative unitary ring, $I \subsetneq R$ be an ideal.
    $$\Theta = \{\text{ideals of $J$ of } R \mid I \subseteq J \neq R\},$$
    $\Theta \neq \varnothing$ and ordered by $\subseteq$.
    If $\Theta_0 \neq \varnothing$, totally ordered subset of $\Theta$ then $\bigcup_{J \Theta_0} J$ is an upper bound of $\Theta_0$ in $\Theta$.
    By Zorn lemma, $\Theta$ admits a maximal element which is necessary a maximal ideal.
\end{remark}

\begin{propositionenv}\label{prop:6.1.16}
    Let $A$ be an principal ideal domain,
    then any $a \in A\setminus A^\times$ can be written as 
    $$a = p_1 \cdots p_n,$$
    where $p_1,\cdots, p_n \in A$ are irreducible.
    Moreover, if
    $$a = p_1 \cdots p_n = q_1 \cdots q_m,$$
    then $n = m$ and there exists $\sigma \in \mathfrak{S}_n$ such that $\forall i \in \{1,\cdots,n\}$,
    $p_i A = q_{\sigma(i)} A.$
\end{propositionenv}
\begin{proofenv}
    \begin{enumerate}
        \item \textit{Existence}
            We only need to consider the case where $a$ is not irreducible, by corollary \ref{cor:6.1.14},
            $aA = (a)$ is strictly contained in a maximal ideal of $A$.
            There exists $a_1 \in A \setminus A^\times$ such that
            $$a = p_1 a_1.$$
            Iterating this strategy, we obtain
            $$a = p_1 \cdots p_i a_i,$$
            where $p_1,\cdots, p_i$ are irreducible in $A$, $a_i \in A \setminus A^\times$.
            This procedure has to terminate, as otherwise we would construct a strictly increasing
            $$aA \subsetneq a_1 A \subsetneq \cdots \subsetneq a_n A \subsetneq \cdots.$$
            Hence, $a = p_1 \cdots p_n$.
        \item \textit{Uniqueness}
            We reason by induction on the minimal number of irreducible factors in the decomposition.
            If $a$ is irreducible, trivial.

            \quad If $a$ is not irreducible,
            assume now that $a = p_1 \cdots p_n$, $n \ge 2$ and $a$ can not be written as a product of $i < n$ irreducible elements.
            By induction hypothesis,
            we assume the uniqueness (up to permutation) for elements that can be written as a product of $i < n$ irreducible elements.
            If $a = q_1 \cdots q_m \in p_n A$,
            there exists $i\{1,\cdots,m\}$ such that $q_i \in p_n A$ (by our definition, we immediately know that a maximal ideal is prime).
            By permutation, without loss of generality, $i = m$, so $q_m A \subseteq p_n A$.
            Since $q_m$ is irreducible, $q_m A$ is maximal too.
            So $q_m A = p_n A$.
            By proposition \ref{prop:6.1.3},
            there exists $\lambda \in A^\times$ such that $q_m = \lambda p_n$.
            So $p_1\cdots p_{n-1} = (\lambda q_1) \cdots q_{m-1}.$
            Note that $q_1 A = (\lambda q_1) A$,
            we obtain the desired result using the induction hypothesis.
    \end{enumerate}
\end{proofenv}

\begin{remark}
    Let $K$ be a field and $M$ be the set of all monic irreducible polynomials in $K[T]$.
    By \ref{6.1.15} and \ref{prop:6.1.16},
    any non-zero $F \in K[T]$ can be decomposed as 
    $$F(T) = c(F) P_1(T) \cdots P_n(T),$$
    where $P_1,\cdots, P_n \in M$.
    It is unique up to a permutation.
    For any $P \in M$, denote by $\ord_P(F)$ the number of occurrence of $P$ in the decomposition of $F$.
    $$F(T) = c(F) \prod_{p \in M} P(T)^{\ord_P(F)}.$$
    If $F,G \in K[T]$ to be non-zero,
    then for any $P \in M$,
    $$\ord(F \cdot G) = \ord_P(F) + \ord_P(G).$$
    $$\ord_P(F + G) \ge \min\{\ord_P(F), \ord_P(G)\}.$$
    By convention, $\ord_P(0) = +\infty$.
    Moreover, $F \mid G$ if and only if
    $$\forall P \in M,\ \ord_P(G) \ge \ord_P(F).$$
    
    \quad Conversely, assume that any $I$ in $R$ is finitely generated.
    Let
    $$I_0 \subseteq I_1 \subseteq \cdots I_n \subseteq I_{n+1} \subseteq \cdots$$
    be an increasing sequence of ideals.
    Then $\dis I \coloneq I_n$ is an ideal in $R$, so $I = (x_1,\cdots,x_p)$.
    Let $n \in \NN$ such that 
    $$I_n \supseteq \{x_1,\cdots x_p\},$$
    hence $(x_1,\cdots,x_p) = I_n = I_{n+1} = \cdots$.
    Then any ideal $I \neq R$ is contained in some maximal ideal of $R$.
\end{remark}

\begin{propositionenv}
    Let $A$ be a commutative unitary ring.
    The set of all ideals of $A$ equipped with the relation of inclusion forms an order complete totally ordered set.
\end{propositionenv}
\begin{proofenv}
    Let $\Theta$ be a family of all ideals of $A$.
    If $\Theta$ is not empty, then $\dis \bigcap_{I \in \Theta} I$ is an ideal in $\Theta$.
    It is the the infimum of $\Theta$.
    Moreover, the ideal generated by $\dis \bigcup_{I \in \Theta } I$ is the supremum of $\Theta$.
    We deduce it by
    $$\sum_{I \in \Theta} I.$$
\end{proofenv}

\begin{definitionenv}
    Let $K$ be a field, $n \in \NN$, and $F_1,\cdots, F_n \in K[T]$.
    We call the \textbf{greatest common divisor} of $F_1,\cdots, F_n$ the unique monic polynomial that generates the ideal
    $$\sum_{i=1}^n F_i(T)K[T]$$
    We denote it $\gcd(F_1,\cdots, F_n)$.
    There exists $G_1,\cdots G_n \in K[T]$ such that
    $$\gcd(F_1,\cdots, F_n)(T) = \sum_{i=1}^n F_i(T)G_i(T).$$
    We call the \textbf{least common multiple} the polynomial that generates
    $$\bigcap_{i=1}^n F_i(T)K[T]$$
    We denote it $\lcm(F_1,\cdots, F_n)$.
\end{definitionenv}

\begin{propositionenv}
    Let $K$ be a field, $n \in \NN_{>0}$, $F_1, \cdots F_n \in K[T]$.
    For any irreducible monic $P \in K[T]$,
    $$\ord_P(\gcd(F_1,\cdots, F_n)) = \min_{i \in \{1,\cdots,n\}} \ord_{P}(F_i),$$
    $$\ord_P(\lcm(F_1,\cdots, F_n)) = \max_{i \in \{1,\cdots,n\}} \ord_{P}(F_i).$$
\end{propositionenv}
\begin{proofenv}
    By definition, for any $i \in \{1,\cdots,n\}$,
    $$\gcd(F_1,\cdots, F_n) \mid F_i,\quad F_i \mid \lcm(F_1,\cdots, F_n).$$
    Hence,
    $$\ord_P\left(\gcd(F_1, \cdots F_n)\right) \le \ord_P(F_i) \quad \text{and} \quad \ord_P\left(\lcm(F_1, \cdots F_n)\right) \ge \ord_P(F_i).$$
    Let $M$ be the set of monic irreducible $P \in K[T]$.
    $$F = \prod_{P \in M} P^{\min\{\ord_P(F_1),\cdots,\ord_P(F_n)\}}.$$
    For any $i \in \{1,\cdots,n\}$, $F \mid F_i$, so $F \mid \gcd(F_1,\cdots,F_n)$,
    $$\min\{\ord_{P}(F_1),\cdots, \ord_{P}(F_n)\} \le \ord_{P}\left(\gcd(F_1,\cdots,F_n)\right).$$
    Let
    $$G = \prod_{P \in M} P^{\max\{\ord_P(F_1),\cdots,\ord_P(F_n)\}}.$$
    For any $i \in \{1,\cdots,n\}$, $G$ is a multiple of $F_i$,
    hence
    $$\lcm(F_1,\cdots,F_n) \mid G.$$
    Therefore,
    $$\max\{\ord_P(F_1),\cdots, \ord_P(F_n)\} \ge \ord_{P}\left(\lcm(F_1,\cdots,F_n)\right).$$
\end{proofenv}


\section{Cayley-Hamilton Theorem}


\begin{definitionenv}\label{def:6.2.1}
    Let $K$ be a commutative unitary ring.
    A $K$-algebra is a $K$-module $A$ with a composition law
    $$\begin{array}{rcl}
        A \times A & \to & A\\
        (a,b) & \mapsto & ab
    \end{array}$$
    such that $A$ with the addition and the composition law forms a unitary ring such that
    $$\forall(\lambda,a,b) \in K \times A \times A,\ \lambda(ab) = \lambda(a)b = a\lambda(b).$$
    Let $A$ be an $K$ algebra and $V$ be a $K$-module.
    We equip $A \otimes V$ with the following structure of left $A$-module
    $$\begin{array}{rcl}
        A \times (A \otimes V) & \longrightarrow & A \otimes V\\
        (a,(b \otimes x)) & \longmapsto & ab \otimes x.
    \end{array}$$
\end{definitionenv}

\begin{propositionenv}\label{prop:6.2.2}
    Let $K$ be a commutative unitary ring,
    $A$ be a $K$ algebra and $V$ be a free $K$-module of finite rank $n$.
    Let $(e_i)_{i=1}^{n}$ be a basis of $V$.
    Then $(1 \otimes e_i)_{i=1}^{n}$ is a basis of $A \otimes V$.
\end{propositionenv}
\begin{proofenv}
    Let $M$ be a left $A$-module.
    One has the following bijections, which are functorial with respect to $M$.
    $$\begin{aligned}
        \Hom_A (A \otimes V, M) \cong & \{f : A \times V \rightarrow M \mid f \text{ is $K$-bilinear, } f (a,x) = a f(1,x)\}\\
        \cong & \Hom_K (V,M) \cong M^{\{e_1,\cdots,e_n\}} \cong M^{\{1 \otimes e_1,\cdots,1 \otimes e_n\}}.
    \end{aligned}$$
    Therefore, $A \otimes V$ is a free $A$-module generated by $\{1 \otimes e_i\}_{i=1}^{n}$.
\end{proofenv}

\begin{definitionenv}\label{def:6.2.3}
    Let $K$ be a commutative unitary ring,
    $V$ be a $K$-module and $\varphi \in \End_{K}(V)$.
    For any polynomial
    $$F(T) = \sum_{i=0}^{d} a_i T^i \in K[T],$$
    we denote by $F(\varphi) \in \End_{K}(V)$ the following $K$-linear endomorphism
    $$F(\varphi) = \sum_{i=0}^{d} a_i \varphi^i,\quad \varphi^0 = \Id_V.$$
    Note that $F \rightarrow F(\varphi)$ defines a homomorphism of $K$-algebras from $K[T]$ to $\End_{K}(V)$.
    Since $K[T]$ is commutative, we obtain that, for any $FG \in K[T]$, 
    $$F(\varphi) \circ G(\varphi) = G(\varphi) \circ F(\varphi).$$
    We equip $V$ with the structure of left $K[T]$-module in the following way
    $$\begin{array}{rcl}
        K[T] \times V & \longrightarrow & V\\
        (F(T),x) & \longmapsto & F(\varphi)(x).
    \end{array}$$
    This mapping is $K$-linear and hence defines a $K$-linear mapping $K[T] \otimes V \rightarrow V$ that sends $F \otimes x$ to $F(\varphi)(x)$.
    Note that if we consider $K[T] \otimes V$ as a $K[T]$-module and consider the above structure of $K[T]$-module on $V$,
    then the $K$-linear mapping $K[T]\otimes V \ni F \otimes x \mapsto F(\varphi)(x) \in V$ is actually $K[T]$-linear.
    Indeed, for any $F,G \in K[T]$, one has
    $$(FG)(\varphi)(x) = (F(\varphi) \circ G(\varphi))(x) = F(\varphi) (G(\varphi))(x).$$
\end{definitionenv}

\begin{lemmaenv}\label{lem:6.2.4}
    Let $K$ be a commutative unitary ring,
    $V$ be a free $K$-module of finite rank $n$ and $\varphi \in \End_{K}(V)$.
    Then one has
    $$(\varphi^a)^\vee = (\varphi^\vee)^a.$$
\end{lemmaenv}
\begin{proofenv}
    Let $(e_i)_{i=1}^{n}$ be a basis of $V$,
    $(e_i^\vee)_{i=1}^{n}$ be its dual basis.
    One has 
    $$e_1^\vee \wedge \varphi^\vee (e_2^\vee) \wedge \cdots \wedge \varphi^\vee (e_n^\vee) = e_1^\vee \wedge (e_2^\vee \circ \varphi) \wedge \cdots \wedge (e_n^\vee \circ \varphi).$$
    Therefore,
    $$\begin{aligned}
        &(e_1^\vee \wedge \varphi^\vee (e_2^\vee) \wedge \cdots \wedge \varphi^\vee (e_n^\vee)) (e_1\wedge \cdots \wedge e_n)\\
        = & \sum_{\sigma \in \mathfrak{S}_n} \sgn(\sigma) e_1^\vee(e_{\sigma(1)})\wedge e_2^\vee(\varphi(e_{\sigma(2)})) \wedge \cdots \wedge e_n^\vee(\varphi(e_{\sigma(n)})).
    \end{aligned}$$
    We obtain
    $$\begin{aligned}
        & (e_1^\vee \wedge \varphi^\vee(e_2^\vee) \wedge \cdots \wedge \varphi^\vee(e_n^\vee))(e_1 \wedge \cdots \wedge e_n)\\
        = & (e_1^\vee \wedge \cdots \wedge e_n^\vee) (e_1 \wedge \varphi(e_2) \wedge \cdots \wedge \varphi(e_n))\\
        = & (e_1^\vee \wedge \cdots \wedge e_n^\vee)(\varphi^a(e_1) \wedge e_2 \wedge \cdots \wedge e_n)\\
        = & \sum_{\sigma \in \mathfrak{S}_n} \sgn(\sigma) e^\vee_{\sigma(1)} (\varphi^a(e_1))e_{\sigma(2)^\vee}(e_2) \cdots e_{\sigma(n)}^\vee(e_n)\\
        = & (e_1^\vee \circ \varphi^a) \wedge e_2^\vee \wedge \cdots \wedge e_n^\vee (e_1 \wedge \cdots \wedge e_n).
    \end{aligned}$$
    We obtain
    $$(\varphi^a)^\vee (e_1^\vee) = (\varphi^\vee)^a (e_1^\vee).$$
\end{proofenv}

\begin{lemmaenv}\label{lem:6.2.5}
    Let $K$ be a commutative unitary ring,
    $V$ be a free $K$-module of finite rank $n$ and $\varphi, \psi \in \End_{K}(V)$.
    If $\varphi^\vee = \psi^\vee$, then $\varphi = \psi$.
\end{lemmaenv}
\begin{proofenv}
    Let $(e_i)_{i=1}^{n}$ be a basis of $V$,
    $(e_i^\vee)_{i=1}^{n}$ be its dual basis.
    By assumption $\varphi^\vee = \psi^\vee$, we have
    $$ e_i^\vee \circ \varphi = e_i^\vee \circ \psi,\; i \in \{1,\cdots,n\}.$$
    For any $x \in V$,
    $$x = \sum_{i=1}^{n} e_i^\vee(x) \cdot e_i.$$
    Therefore, for any $y \in V$,
    $$\varphi(y) = \sum_{i=1}^{n} (e_i^\vee \circ \varphi) (y) = \sum_{i = 1}^{n} (e_i^\vee \circ \psi) (y) = \psi(y).$$
\end{proofenv}

\begin{propositionenv}\label{prop:6.2.6}
    Let $K$ be a commutative unitary ring,
    $V$ be a free $K$-module of finite rank $n$ and $\varphi \in \End_{K}(V)$.
    Then
    $$\varphi \circ \varphi^a = \det(\varphi) \Id_V.$$
\end{propositionenv}
\begin{proofenv}
    Applying the theorem \ref{thm:4.7.1} to $\varphi^\vee$, by lemma \ref{lem:6.2.4} and proposition \ref{thm:LinearMap.det_dualMap}, 
    $$(\varphi \circ \varphi^a)^\vee = (\varphi^a)^\vee \circ \varphi^\vee = \det(\varphi^\vee) \Id_{V^\vee} = \det(\varphi) \Id_{V^\vee}.$$
    Hence, by lemma \ref{lem:6.2.5}, we obtain the claimed equality.
\end{proofenv}

\begin{theoremenv}[Cayley-Hamilton]
    Let $K$ be a commutative unitary ring,
    $V$ be a free $K$-module of finite rank $n$ and $\varphi \in \End_{K}(V)$.
    Let $T \cdot \Id_{V} - \varphi$ denote the $K[T]$-linear endomorphism
    $$T \otimes \Id_{V} - 1 \otimes \varphi$$
    of $K[T] \otimes V$ that sends $F(T) \otimes x$ to 
    $$TF(T) \otimes x - F(T) \otimes \varphi(x),$$
    and let
    $$P_\varphi(T) = \det(T \cdot \Id_{V} - \varphi).$$
    Then
    $$P_\varphi(\varphi) = 0$$
    in $\End_{K}(V)$.
\end{theoremenv}
\begin{proofenv}
    For any element of the form $F(T) \otimes x \in K[T] \otimes V$,
    the image of 
    $$(T\cdot \Id_{V} - \varphi)(F(T) \otimes x)$$
    in $V$ by the $K[T]$-linear mapping $K[T] \otimes V \rightarrow V$
    described in definition \ref{def:6.2.3} vanishes, cause
    $$\varphi \circ F(\varphi) = F(\varphi) \circ \varphi.$$
    Indeed,
    $$(T\Id_{V} - \varphi)(F \otimes x) = T (F \otimes x) - F \otimes \varphi(x),$$
    $$\mu(T \Id_{V} - \varphi)(F \otimes x) = \varphi F (\varphi)(x) - F(\varphi)(\varphi(x)).$$
    By linearity, we deduce that for any $\xi \in K[T] \otimes V$,
    the image of $(T \Id_{V} - \varphi) (\xi)$ by the $K[T]$-linear mapping
    $$\begin{array}{rrcl}
        \mu : & K[T] \otimes V & \longrightarrow & V \\
        & F(T) \otimes x & \longmapsto & F(\varphi)(x)
    \end{array}$$
    vanishes.
    By proposition \ref{prop:6.2.6}, 
    $$(T - \varphi) \circ (T - \varphi)^a = P_{\varphi}(T) \Id_{K[T] \otimes V}.$$
    Let $x \in V$, consider $1 \otimes x$, then
    $$P_{\varphi} (T) \otimes x = P_{\varphi}(T)(1 \otimes x) = (T\Id_{V} - \varphi)(\xi),$$
    where $\xi = (T\Id_{V} - \varphi)^a(1 \otimes x)$.
    Taking the image of this under $\mu$, we get
    $$P_{\varphi}(\varphi)(x) = 0.$$
\end{proofenv}

\begin{definitionenv}
    Let $K$ be a commutative unitary ring,
    $V$ be a free $K$-module of finite rank $n$ and $\varphi \in \End_{K}(V)$.
    The polynomial $P_\varphi(T) = \det(T \cdot \Id_{V} - \varphi)$ is called the \textbf{characteristic polynomial} of $\varphi$.

    \quad Let $(e_i)_{i=1}^n$ be a basis of $V$ over $K$.
    Then $(1 \otimes e_1)_{i=1}^{n}$ forms a basis of the $K[T]$-module $K[T] \otimes V$.
    By abuse of notation, for any $x \in V$,
    we still use the notation $x$ to denote the element $1 \otimes x$ in $K[T] \otimes V$.
    We therefore identify $V$ as a sub-$K$-module of $K[T] \otimes V$ with this notation
    $$(T\Id_{V} - \varphi)(e_i) = Te_i - \varphi(e_i).$$
    Therefore,
    $$(T\Id_{V} - \varphi)(e_1) \wedge \cdots \wedge (T\Id_{V} - \varphi)(e_n) = (Te_1 - \varphi(e_1)) \wedge \cdots \wedge (Te_n - \varphi(e_n)).$$
    This can be written as 
    $$\begin{aligned}
        &T^n(e_1 \wedge \cdots \wedge e_n)\\
        & - T^{n-1} \sum_{i=1}^{n}(e_1 \wedge \cdots \wedge e_{i-1} \wedge \varphi(e_i) \wedge e_{i+1} \wedge \cdots \wedge e_n) + \cdots \\
        &+ (-1)^n \varphi(e_1) \wedge \cdots \wedge \varphi(e_n).
    \end{aligned}$$
    We see that $P_{\varphi}$ is a monic polynomial of degree $n$,
    and the coefficient of the polynomial is $(-1)^n \det(\varphi)$.
    We denote by $\tr(\varphi)$ the opposite of the coefficient of $T^{n-1}$ in $P_{\varphi}$,
    namely,
    $$(\tr (\varphi))(e_1 \wedge \cdots \wedge e_n) = \sum_{i=1}^{n} (e_1 \wedge \cdots \wedge e_{i-1} \wedge \varphi(e_i) \wedge e_{i+1} \wedge \cdots \wedge e_n).$$
    We can write $P_{\varphi} (T)$ as
    $$P_{\varphi}(T) = T^n - s_1 (\varphi) T^{n-1} + \cdots + (-1)^{n} s_n (\varphi).$$ 
    In particular, $\tr(\varphi) = s_1 (\varphi)$ and $\det(\varphi) = s_n(\varphi).$
\end{definitionenv}

\begin{remark}
    Assume that
    $$P_{\varphi} (T) = (T - \lambda_1) \cdots (T - \lambda_n),$$
    then
    $$s_{l} = \sum_{1 \le i_1 \le \cdots \le i_{l} \le n} \lambda_{i_1} \cdots \lambda_{i_{l}}.$$
    By definition, for any $(x_{1},\cdots,x_{n}) \in V^n$,
    $$(Tx_1 + \varphi(x_1)) \wedge \cdots \wedge (Tx_n + \varphi(x_n)) = (T^n + s_1(\varphi) T^{n-1} + \cdots + s_d(\varphi)) x_1 \wedge \cdots \wedge x_{n}.$$
    Hence,
    $$s_l (\varphi)(x_1 \wedge \cdots \wedge x_n) = \sum_{\substack{I = (i_1,\cdots,i_{l})\\1\le i_1 \le \cdots \le i_{l} \le n}} z_1^I \wedge \cdots \wedge z_n^I,$$
    where 
    $$z_j^I = 
    \begin{cases}
        x_j, \ &j \notin I\\
        \varphi(x_j), \ &j \in I
    \end{cases}.$$
    Let $F_{l} : \End_{K}(V)^l \rightarrow K$ be the multilinear mapping
    $$F_{l}(\varphi_1, \cdots, \varphi_l) x_1 \wedge \cdots x_d = \sum_{\substack{I = (i_1,\cdots,i_{l})\\1\le i_1 \le \cdots \le i_{l} \le n}} u_1^I \wedge \cdots \wedge u_d^I,$$
    where 
    $$u_j^I = 
    \begin{cases}
        x_j \ &j \notin I\\
        \varphi_k(x_j) \ &j \in I
    \end{cases}.$$
    $F_l$ is multilinear form and $s_l (\varphi) = F_{l}(\varphi,\cdots,\varphi)$.
    Hence $s_l$ is a homogenous polynomial of degree $l$ on $\End_{K}(V)$.
\end{remark}


\section{Reduction of Linear Endomorphisms}


In this section, we fix a field $K$.

\begin{definitionenv}\label{def:6.3.1}
    Let $V$ be a finite dimensional vector space over $K$.
    If $\varphi$ is a $K$-linear endomorphism of $V$,
    we let $\lambda_{\varphi}$ to be the homomorphism of $K$-algebras
    $$\begin{array}{rrcl}
        \lambda_{\varphi} : & K[T] & \longrightarrow & \End_{K}(V)\\
        & F(T) & \longmapsto & F(\varphi).
    \end{array}$$
    By Cayley-Hamilton theorem, $P_{\varphi} \in \ker(\lambda_{\varphi})$.
    As a consequence, $\ker (\lambda_{\varphi})$ is a non-zero ideal in $K[T]$ (a principal ideal domain).
    We denote by $Q_{\varphi}(T)$ the unique monic polynomial that generates $\ker(\lambda_{\varphi})$.
    
    \quad Note that $F \in K[T]$ satisfies $F(\varphi) = 0$ if and only if $Q_{\varphi} \mid F$.
    We call $Q_\varphi$ the \textbf{minimal polynomial} of $\varphi$.
\end{definitionenv}

\begin{propositionenv}\label{prop:6.3.2}
    Let $V$ be a finite dimensional vector space over $K$,
    $\varphi \in \End_{K}(V)$,
    and $W$ be a vector subspace of $V$ such that $\varphi(W) \subseteq W$.
    Let $\psi : W \rightarrow W$ be the restriction of $V$ to $W$.
    Then $Q_{\psi} \mid Q_{\varphi}$.
\end{propositionenv}
\begin{proofenv}
    $Q_{\varphi}(\varphi) = 0$, hence, for any $x \in W$,
    $$Q_{\varphi}(\psi)(x) = Q_{\varphi}(\varphi)(x) = 0.$$
    Hence, $Q_{\varphi}(\psi) = 0$, then $Q_{\psi} \mid Q_{\varphi}$.
\end{proofenv}

\begin{definitionenv}\label{def:6.3.3}
    Let $V$ be a vector space over $K$,
    $I$ be a set and $(V_i)_{i \in I}$ be a family of vector subspaces of $V$.
    Recall that, in the case where the $K$-linear mapping
    $$ \bigoplus_{i \in I} V_i \longrightarrow V, \quad (x_i)_{i \in I} \longmapsto \sum_{i \in I} x_i$$
    is an isomorphism of $K$-vector spaces, we say that $V$ is a direct sum of $(V_i)_{i \in I}$.
    
    \quad Note that if for any $i \in I$, $(e_{i,j})_{j \in J_i}$ is a basis of $V_i$, then
    $$e_{i,j}, \quad i \in I, j \in J_i$$
    forms a basis of $V$ over $K$.
\end{definitionenv}

\begin{propositionenv}\label{prop:6.3.4}
    Let $V$ be a finite dimensional vector space over $K$,
    $\varphi \in \End_{K}(V)$,
    $V_1,\cdots,V_n$ be vector subspaces of $V$.
    We assume that, for any $i \in \{1,\cdots,n\}$, $\varphi(V_i) \subseteq V_{i}$.
    Let $\varphi_i:V_i \rightarrow V_i$ be the restriction of $\varphi$ to $V_i$.
    \begin{enumerate}
        \item If $V = V_1 + \cdots + V_n$, then $Q_{\varphi} (T) = \lcm(Q_{\varphi_1}(T), \cdots, Q_{\varphi_n}(T))$.
        \item If $V$ is the direct sum of $V_1,\cdots,V_n$, then 
            $$P_{\varphi}(T) = P_{\varphi_1}(T) \cdots P_{\varphi_n}(T).$$
    \end{enumerate}
\end{propositionenv}
\begin{proofenv}
    \quad
    \begin{enumerate}
        \item Let $G = \lcm(Q_{\varphi_1}, \cdots, Q_{\varphi_n})$.
            For any $i \in \{1,\cdots,n\}$,
            let $F_i \in K[T]$ be such that $G = F_{i} Q_{\varphi_i}$.
            Let $x \in V$, $x = x_1 + \cdots + x_n$, $x_i \in V_i$, $i \in \{1,\cdots,n\}$.
            One has
            $$G(\varphi)(x) = \sum_{i=1}^{n} G(\varphi)(x_i) = \sum_{i=1}^{n} G(\varphi_i)(x_i) = \sum_{i=1}^{n} F_i (\varphi_i) Q_{\varphi_i}(\varphi_i)(x).$$
            Hence, $G$ is a vanishing polynomial of $\varphi$, so $Q_\varphi \mid G$.
            Conversely, by proposition \ref{prop:6.3.2}, one has that $Q_{\varphi_i} \mid Q_{\varphi}$ for any $i \in \{1,\cdots,n\}$.
            Hence, $G \mid Q_{\varphi}$.
        \item For any $i \in \{1,\cdots,n\}$, let $(e_{i,j})_{j=1}^{d_i}$ be a basis of $V_i$, $d_i = \dim(V_i)$.
            One has
            $$(Te_{i,1} - \varphi(e_{i,1})) \wedge \cdots \wedge (Te_{i,d_i} - \varphi(e_{i,d_i})) = P_{\varphi_i} (T) e_{i,1} \wedge \cdots \wedge e_{i,d_i}.$$
            $$P_{\varphi}(T) = P_{\varphi_{1}}(T) \cdots P_{\varphi_{n}}(T).$$
    \end{enumerate}
\end{proofenv}

\begin{theoremenv}
    Let $V$ be a finite dimensional vector space over $K$,
    $\varphi \in \End_{K}(V)$,
    let $F_1,\cdots,F_n$ be non-zero polynomials in $K[T]$ and $F = F_1 \cdots F_n$.
    Assume that $F_1,\cdots F_n$ are pairwise coprime,
    that is for any $i \neq j \in \{1,\cdots,n\}$, $F_i,F_j$ do not have any non-trivial common irreducible factor.
    Let $F = \prod_{i=1}^{n} F_i$.
    Then $\ker(F(\varphi))$ is the direct sum of
    $$\ker(F_1(\varphi)), \cdots , \ker(F_n(\varphi)).$$
\end{theoremenv}
\begin{proofenv}
    For any $i \in \{1,\cdots,n\}$,
    $$G_{i}(T) = F_1(T) \cdots F_{i-1}(T) F_{i+1}(T) \cdots F_n(T).$$
    $$\gcd(G_{1}(T), \cdots, G_{n}(T)) = 1.$$
    Hence there exist polynomials $A_1,\cdots,A_n$ in $K[T]$ such that
    $$1 = A_1(T) G_1(T) + \cdots + A_n(T)G_n(T).$$
    So $\Id_{V} = A_1(\varphi) \circ G_{1}(\varphi) + \cdots + A_n(\varphi) \circ G_{n}(\varphi).$
    Namely, $\forall x \in V$,
    $$x = A_1(\varphi) \left(G_1(\varphi)(x)\right) + \cdots + A_n(\varphi) \left(G_n(\varphi)(x)\right).$$
    If $x \in \ker\left(F(\varphi)\right)$,
    then for any $i \in \{1,\cdots,n\}$, 
    $$F_{i}(\varphi)\left(A_{i}(\varphi)\left(G_i(\varphi)(x)\right)\right) = A_i(\varphi)\left(F_i(\varphi)\left(G_i(\varphi)(x)\right)\right) = 0.$$
    So
    $$A_i(\varphi)\left(G_{i}(\varphi)(x)\right) \in \ker\left(F_{i}(\varphi)\right).$$
    Suppose that $(x_1,\cdots,x_n)\in \ker(F_1(\varphi)) \times \cdots \times \ker(F_{n}(\varphi))$ such that $x = x_1 + \cdots + x_n$.
    Then
    $$\begin{aligned}
        & \sum_{j=1}^{n}A_i(\varphi)\left(G_i(\varphi)(x_j)\right)\\
        = & A_i(\varphi)\left(G_i(\varphi)(x_i)\right)\\
        = & \sum_{j=1}^{n} A_{j}(\varphi)\left(G_j(\varphi)(x_i)\right)\\
        =& x_i.
    \end{aligned}$$
\end{proofenv}

\begin{remark}
    \quad
    \begin{enumerate}
        \item If $x \in \ker(\varphi)$, then
            $$F(\varphi)\left(\varphi(x)\right) = \varphi\left(F(\varphi)(x)\right) = 0.$$
            So $\varphi (x) \in \ker(F(\varphi))$.
        \item Suppose that $F(\varphi) = 0$,
            then $V$ is the direct sum of $\ker(F_i(\varphi))$, $i \in \{1,\cdots,n\}$.
    \end{enumerate}
\end{remark}

\begin{propositionenv}
    Let $K$ be a field and $V$ be a finite dimensional vector space over $K$.
    Let $\varphi \in \End_{K}(V)$ and $P_{\varphi}$ and $Q_{\varphi}$ be the characteristic polynomial and the minimal polynomial of $\varphi$, respectively.
    Then $Q_{\varphi} \mid P_{\varphi} \mid Q_{\varphi}^{\dim(V)}$.
    In particular, $P_{\varphi}$ and $Q_{\varphi}$ have the same irreducible factors.
\end{propositionenv}
\begin{proofenv}
    Suppose that $Q_{\varphi}(T) = T^n + a_1 T^{n-1} \cdots + \cdots + a_n$.
    $$\begin{aligned}
        & Q_{\varphi}(T) \Id_{K[T]\otimes V} = Q_{\varphi}(T) \otimes \Id_{V} - 1 \otimes Q_{\varphi}(\varphi)\\
        = & \left(T^n + a_1 T^{n-1} + \cdots + a_n\right) \otimes \Id_{V} - 1 \otimes \left(\varphi^n + a_1 \varphi^{n-1} + \cdots + a_n\right)\\
        = & \left(T^n \otimes \Id_{V} - 1 \otimes \varphi^n\right) + \sum_{i=1}^{n}a_i \left(T^{n-i}\otimes\Id_{V} - 1 \otimes \varphi^{n-i}\right)\\
        = & \left(T \otimes \Id_{V} - 1 \otimes \varphi\right) \circ \left(\sum_{i=0}^{n-1}a_i \sum_{j=0}^{n-i-1}\left(T^j \otimes \varphi^{n-i-1-j}\right)\right).
    \end{aligned}$$
    Taking the determinant, we get
    $$Q_{\varphi}(T)^{\dim(V)} = P_\varphi(T) \cdot \det\left(\sum_{i=0}^{n-1}a_i \sum_{j=0}^{n-i-1}\left(T^j \otimes \varphi^{n-i-1-j}\right)\right).$$
\end{proofenv}

\begin{definitionenv}
    Let $K$ be a field and $V$ be a finite dimensional vector space over $K$.
    Let $\varphi \in \End_{K}(V)$.
    If there is a basis $(e_i)_{i=1}^n$ which is composed of eigenvectors of $\varphi$, then we say that $\varphi$ is
    \textbf{diagonalizable} (or semi-simple).
\end{definitionenv}

\begin{propositionenv}
    Let $K$ be a field and $V$ be a finite dimensional vector space over $K$,
    and $\varphi \in \End_{K}(V)$
    Let $\lambda \in K$.
    The following statements are equivalent:
    \begin{enumerate}[label = (\arabic*)]
        \item $\lambda$ is an eigenvalue of $\varphi$.
        \item $\lambda$ is a root of the minimal polynomial $Q_{\varphi}$.
        \item $\lambda$ is a root of the characteristic polynomial $P_{\varphi}$.
    \end{enumerate}
\end{propositionenv}

\begin{proofenv}
    (1) $\Rightarrow$ (2). Let $W \coloneq \ker(\varphi - \lambda \Id_{V}) \neq 0$,
    $Q_{\varphi}(T) = T - \lambda$ since $\varphi|_{W} = \lambda \Id_{W}$.
    Since $Q_{\varphi}|_{W} \mid Q_{\varphi}$, $Q_{\varphi}(\lambda) = 0$.

    \quad
    (2) $\Rightarrow$ (1). Suppose that $Q_{\varphi}(T) = (T - \lambda)F(T)$.
    $\ker(F(\varphi)) \neq V$, so there exists $y \in V$, such that $x \coloneq F(\varphi)(y) \neq 0$.
    $$(\varphi - \lambda \Id_{V})(x) = (\varphi - \lambda \Id_{V})(F(\varphi)(y)) = Q_{\varphi}(\varphi)(y).$$
\end{proofenv}

\begin{theoremenv}
    Let $K$ be a field, $V$ be a finite-dimensional vector space over $K$,
    and $\varphi \in \End_{K}(V)$.
    Then $\varphi$ is diagonalizable if and only if all roots of the minimal polynomial of $Q_\varphi$ are simple,
    namely, $Q_{\varphi}$ can be written into the form
    $$Q_{\varphi}(T) = (T - \lambda_1)(T - \lambda_2)\cdots(T - \lambda_n),$$
    where $\lambda_1, \lambda_2, \cdots, \lambda_n$ are distinct elements of $K$.
\end{theoremenv}
\begin{proofenv}
    Suppose that $(e_i)_{i=1}^d$ is a basis of $V$ which is composed of eigenvectors of $\varphi$ with eigenvalues $\mu_1,\cdots,\mu_d$.
    For any $i \in \{1,\cdots,n\}$, let $W_i = K e_i$, $\varphi_i = \varphi|_{W_i} = \mu_i \Id_{W_i}$.
    So $Q_{\varphi_i} = T - \mu_i$,
    $$Q_{\varphi} = \lcm(Q_{\varphi_1}, \cdots, Q_{\varphi_d}) = \lcm(T-\mu_1, \cdots, T-\mu_d).$$
    Suppose that
    $$Q_{\varphi}(T) = (T - \lambda_1)(T - \lambda_2)\cdots(T - \lambda_n),$$
    where $\lambda_1, \lambda_2, \cdots, \lambda_n$ are distinct.
    Hence $T-\lambda_1,\cdots,T-\lambda_n$ are pairwise coprime.
    So $V$ is the direct sum of eigenspaces $\ker(\varphi - \lambda_i \Id_{V}),\cdots,\ker(\varphi - \lambda_n \Id_{V})$.
    Taking a basis $e_i$ of $\ker(\varphi - \lambda_i \Id_{V})$,
    then $\dis \bigcup_{i=1}^{n}e_i$ is a basis of $V$,
    which is composed of eigenvectors.
\end{proofenv}

\begin{remark}
    We denote by $\Sp(\varphi)$ the set of $\lambda \in K$ such that $\varphi - \lambda \Id_{V}$ is not injective (namely $\lambda$ is an eigenvalue of $\varphi$).
    If $\varphi$ is diagonalizable, then
    $$Q_{\varphi} (T) = \prod_{\lambda \in \Sp(\varphi)} (T - \lambda)$$
\end{remark}


\section{Dunford Decomposition}


We fix a field $K$.

\begin{definitionenv}
    Let $V$ be a finite-dimensional vector space over $K$ and $\varphi \in \End_{K}(V)$.
    Suppose that there exists $k \in \NN_{\ge 1}$ such that $\varphi^k = 0$, then we say that $\varphi$ is \textbf{nilpotent}.
\end{definitionenv}

\begin{remark}
    Suppose that $\varphi$ is nilpotent.
    Then $Q_{\varphi}$ of the form $T^n$.
    So
    $$Q_\varphi(T) = T^n \mid P_{\varphi} \mid Q_{\varphi}^{\dim(V)} = T^{n \dim(V)}.$$
    So $P_{\varphi} = T^{\dim(V)}$.
\end{remark}

\begin{theoremenv}[Dunford decomposition, semi-simple-nilpotent decomposition]
    Let $V$ be a finite-dimensional vector space over $K$ and $\varphi \in \End_{K}(V)$.
    Suppose that the characteristic polynomial $P_{\varphi}$ is of the form
    $$P_{\varphi}(T) = (T - \lambda_1)^{\alpha_1}\cdots(T - \lambda_r)^{\alpha_r},$$
    where $\lambda_1,\cdots,\lambda_r$ are distinct and $\alpha_1,\cdots,\alpha_r$ are positive integers.
    Then there exists $\mu$ and $\nu$ in $K[T]$ such that
    \begin{enumerate}[label = (\arabic*)]
        \item $s \coloneq \mu(\varphi)$ is diagonalizable.
        \item $n \coloneq \nu(\varphi)$ is nilpotent.
        \item $\varphi = s + n$.
    \end{enumerate}
\end{theoremenv}
\begin{remark}
    \quad
    \begin{enumerate}
        \item There is a unique pair $(s,n)$ in $\End_{K}(V)$ such that
        \begin{itemize}
            \item $s$ is diagonalizable
            \item $n$ is nilpotent
            \item $\varphi = s + n$
            \item $s \circ n = n \circ s$
        \end{itemize}
    \end{enumerate}
\end{remark}
\begin{proofenv}
    Let $G_i = P_{\varphi}(T)/(T - \lambda_i)^{\alpha_i}$,
    there exists $(A_1,\cdots,A_r) \in K[T]^r$ such that
    $$1 = \sum_{i=1}^{r} A_i G_i.$$
    Let $\pi_i = A_i (\varphi) \circ G_i (\varphi)$, $\pi_i (V) = V_i$.
    Let $s = \sum_{i=1}^{n} \lambda_i \pi_i$ ($\mu = \sum_{i=1}^{n}\lambda_i A_i G_i$), $\nu = T - \mu(T)$, $n = \varphi - s$.
    One has $\varphi = \varphi \circ \Id_{V} = \sum_{i=1}^{r}\varphi \circ \pi_i$,
    $$n = \varphi - s = \sum_{i=1}^{r}(\varphi - \lambda_i \Id_{V}) \circ \pi_{i}.$$
    So $n^d = \sum_{i=1}^{r} \left(\varphi - \lambda_i \Id_{V}\right)^d \circ \pi_{i} = 0$.
    In fact, if we let $V_i \coloneq \pi_i (V) = \ker\left((\varphi - \lambda_i \Id_{V})^\alpha_i\right)$.

    \quad If $\varphi$ has another decomposition $\varphi = s' + n'$, where $s'$ is semi-simple, $n'$ is nilpotent, and $s' \circ n' = n' \circ s'$.
    Then $s' \circ \varphi = \varphi \circ s'$, $n' \circ \varphi = \varphi \circ n'$.
    Let $x$ be an eigenvector of $s'$ and $\lambda$ be the corresponding eigenvalue.
    Then
    $$n'(x) = \varphi(x) - s' (x) = \varphi (x) - \lambda x = (\varphi - \lambda \Id_{V})(x).$$
    Let $d = \dim(V)$,
    $$0 = (n')^d(x) = (\varphi - \lambda \Id_{V})^d(x).$$
    Hence $x \in \ker((\varphi - \lambda \Id_{V})^d)$, which implies that $(-1)^d P_{\varphi}(\lambda) = \det(\varphi - \lambda \Id_{V}) = 0$.
    Thus $\lambda \in \{\lambda_1,\cdots,\lambda_r\}$.
    For any $i \in \{1,\cdots,r\}$, let
    $$V_i = \ker\left((\varphi - \lambda_i \Id_{V})^{\alpha_i}\right) \subseteq \ker\left((\varphi - \lambda_i \Id_{V})^{d}\right) = V_i'.$$
    $V$ is the direct sum of $V_1,\cdots,V_r$, is also the direct sum of $V_1',\cdots,V_r'$.
    $V_i = V_i'$.
    Since $s'$ is semi-simple, $V$ is the direct sum of $E_{\lambda_i}(s') = \{x \in V \mid s'(x) = \lambda_i x\} \subseteq V_i$.
    So $V_i = E_{\lambda_i}(s')$.
    Hence $s' = s$.
\end{proofenv}


\section{Jordan-Chevalley Decomposition}


We fix a field $K$.

\begin{definitionenv}
    Let $V$ be a finite dimensional vector space over $K$ and $\varphi \in \End_{K}(V)$.
    We say that $\varphi$ is a \textbf{Jordan endomorphism} (Jordan bloc) if both the characteristic polynomial and the minimal polynomial of $\varphi$ are of the form $(T - \lambda)^d$,
    where $\lambda \in K$ and $d = \dim(V)$.
\end{definitionenv}

\begin{lemmaenv}
    Let $V$ be a finite dimensional vector space over $K$,
    $\varphi \in \End_{K}(V)$, $n \in \NN_{\ge 1}$, and $x \in V$.
    Suppose that $\varphi(x)^n = 0$ and $\varphi^{n-1}(x) \neq 0$.
    Then $x,\varphi(x), \cdots , \varphi^{n-1}(x)$ form a linear independent family.
\end{lemmaenv}

\begin{proofenv}
    Let $(a_1,\cdots,a_n) \in K^n$ such that
    $$a_0 x +a_1 \varphi(x) + \cdots + a_{n-1}\varphi^{n-1}(x) = 0.$$
    Then
    $$\varphi^{n-1}\left(a_0 x +a_1 \varphi(x) + \cdots + a_{n-1}\varphi^{n-1}(x)\right) = 0.$$
    So $a_0 = 0$,
    $$a_1 \varphi(x) + \cdots + a_{n-1}\varphi^{n-1}(x) = 0.$$
    By induction, $a_0 = a_1 = \cdots = a_{n-1} = 0$.
\end{proofenv}

\begin{remark}
    Let $V$ be a finite dimensional vector space over $K$ and $\varphi \in \End_{K}(V)$ be a Jordan endomorphism with $P_{\varphi}(T) = (T - \lambda)^d$ with $\lambda \in K$, $d = \dim(V)$.
    Since $P_{\varphi}$ is also the minimal polynomial,
    $$(\varphi - \lambda \Id_{V})^{d - 1} \neq 0.$$
    Pick $x \in V \setminus \ker((\varphi - \lambda \Id_{V})^{d - 1})$.
    By the above lemma, if we denote by $\psi = \varphi - \lambda \Id_{V}$,
    then $x,\psi(x),\cdots,\psi^{d-1}(x)$ form a basis of $V$.
    $$\varphi(\psi^i(x)) = (\lambda\Id_{V} + \psi)(\psi^i(x)) = \psi^{i+1}(x) + \lambda \psi^i(x).$$
    In matrix form,
    $$\begin{pmatrix}
        \varphi(x)\\
        \varphi(\psi(x))\\
        \vdots\\
        \varphi(\psi^{d-1}(x))
    \end{pmatrix} = 
    \begin{pmatrix}
        \lambda x + \psi(x)\\
        \lambda \psi(x) + \psi^2(x)\\
        \vdots\\
        \lambda \psi^{d-1}(x)
    \end{pmatrix}
    =
    \begin{pmatrix}
        \lambda & 1 & 0 & \cdots & 0\\
        0 & \lambda & 1 & \cdots & 0\\
        \vdots & \vdots & \vdots & \ddots & \vdots\\
        0 & 0 & 0 & \cdots & \lambda
    \end{pmatrix}
    \begin{pmatrix}
        x\\
        \psi(x)\\
        \vdots\\
        \psi^{d-1}(x)
    \end{pmatrix}.$$
\end{remark}

\begin{theoremenv}
    Let $V$ be a finite dimensional vector space over $K$ and $\varphi \in \End_{K}(V)$.
    Suppose that the characteristic polynomial of $V$ can be split into a product of polynomial of degree $1$.
    There exist vector subspaces $V_1,\cdots, V_l$ of $V$ such that
    \begin{enumerate}[label = (\arabic*)]
        \item For any $i \in \{1,\cdots,n\}$, $\varphi(V_i) \subseteq V_i$, $\varphi_i \coloneq \varphi|_{V_i} \in \End_{K}(V_i)$ is a Jordan endomorphism.
        \item $V$ is the direct sum of $V_1,\cdots, V_l$.
    \end{enumerate}
\end{theoremenv}

\begin{proofenv}
    We may assume without loss of generality that the minimal polynomial of $\varphi$ is of the form $(T - \lambda)^\alpha$ with $\lambda \in K$, $\alpha \in \NN_{\ge 1}$.
    Let $d = \dim(V)$, $\psi = \varphi - \lambda \Id_{V}$.
    We reason by induction on $n$.

    \quad If $d = 0$, $P_{\varphi} = Q_{\varphi} = 1$.
    If $d = 1$, $P_{\varphi} = Q_{\varphi} = T - \lambda$. In these two cases $\varphi$ is Jordan.

    \quad We suppose that $\varphi$ is NOT Jordan ($d > \alpha$).
    Let $x \in V \setminus \ker(\psi^{\alpha - 1})$ and 
    $$W = \Span_{K}(\{x,\psi(x),\cdots,\psi^{\alpha - 1}(x)\}) \subsetneq V.$$
    Moreover, $(x,\psi(x),\cdots,\psi^{d-1}(x))$ forms a basis of $W$.
    Let $\varphi':W\rightarrow W$ be the restriction of $\varphi$.
    Then the minimal polynomial of $\varphi'$ is $(T - \lambda)^\alpha$.
    Since $\alpha = \dim(W)$, $\varphi'$ is Jordan.
    Let $f: V \rightarrow K$ be a linear form such that
    $$f(\psi^{\alpha - 1}(\alpha)) \neq 0.$$
    Let $H$ be the vector subspace of $V^\vee$ generated by 
    $$\{f,f\circ \psi , f \circ \psi^2 ,\cdots, f \circ \psi^{\alpha - 1}\} = \{f, \psi^\vee (f), (\psi^\vee)^2 (f),\cdots, (\psi^\vee)^{\alpha - 1}(f)\}.$$
    Note that $(\psi^\vee)^\alpha = (\psi^\alpha)^\vee = 0$ and 
    $$(\psi^\vee)^{\alpha - 1}(f)(x) = f(\psi^{\alpha-1}(x)) \neq 0.$$
    By the previous lemma, $f,\psi^\vee(f),\cdots,(\psi^\vee)^{\alpha - 1}(f)$ are $K$-linearly independent,
    hence forms a basis of $H$.
    Let
    $$U = \{y \in V \mid \forall h \in H, h(y) = 0\},$$
    then $\dim(U) = d - \alpha$.
    $$\begin{array}{rrcl}
        \theta : & V & \longrightarrow & K^\alpha\\
        & y & \longmapsto & (f(y), \psi^\vee(f)(y), \cdots, (\psi^\vee)^{\alpha - 1}(f)(y)).
    \end{array}$$
    is surjective since $f, \psi^\vee(f), \cdots, (\psi^\vee)^{\alpha - 1}(f)$ are linearly independent.
    $U = \ker(\theta)$.
    So $\dim(V) = \dim(U) + \alpha$.
    Let $(b_1,\cdots,b_\alpha) \in K^\alpha$, let $i$ be the smallest index in $\{1,\cdots,\alpha\}$ such that $b_i \neq 0$.
    Let $y = b_1 x + b_2 \psi(x) + \cdots + b_{\alpha}\psi^{\alpha-1}(x) \in W$.
    $$f \circ \psi^{\alpha - 1}(y) = b_i (f \circ \psi^{\alpha - 1})(x) \neq 0.$$
    So $y \notin U$, which means $U \cap W = \{0\}$.
    Since $\dim(U) + \dim (W) = \dim (V)$, $V$ is a direct sum of $U$ and $W$.

    If $y \in U$, then $f(y) = f(\psi(y)) = \cdots = f (\psi^{\alpha - 1}(y)) = 0$.
    Hence $f(\psi(y)) = \cdots = f (\psi^{\alpha}(y)) = 0$.
    So $\psi(y) \in U$.
    Hence $\varphi(y) = \lambda y + \psi(y) \in U$.
    Since $U \subseteq V$, one has
    $$Q_{\varphi|_{U}} \mid Q_{\varphi}.$$
    By the induction hypothesis, $U$ can be decomposed into a direct sum of Jordan basis.
\end{proofenv}

\begin{definitionenv}
    The vector space in the above theorem are called Jordan blocs.
\end{definitionenv}

\begin{remark}
    Under the hypothesis of the theorem $\varphi$ has a Dunford decomposition $\varphi = s + n$
    \begin{enumerate}[label = (\arabic*)]
        \item $\dim(\ker(n))$ is the number of Jordan blocs.
        \item $\dim(\ker(n^i)) - \dim(\ker(n^{i-1}))$ is the number of Jordan blocs of dimension greater than $i$.
    \end{enumerate}
\end{remark}


\section{Behavior under Conjugation}


We fix a field $K$ and a finite dimensional vector space $V$ over $K$.
Let $\varphi \in \End_{K}(V)$ and $g \in \GL(V)$, we call conjugate of $\varphi$ by $g$ the element
$$g \circ \varphi \circ g^{-1} \in \End{K}(V).$$
$$\begin{array}{rcl}
    \GL(V) \times \End_{K}(V) & \longrightarrow & \End_{K}(V)\\
    (g,\varphi) & \longmapsto & g \circ g \circ \varphi \circ g^{-1}
\end{array}$$
is a right action.

\begin{theoremenv}
    $$P_{g \circ \varphi \circ g^{-1}} = P_{\varphi}, \quad Q_{g \circ \varphi \circ g^{-1}} = Q_{\varphi}.$$
\end{theoremenv}
\begin{proofenv}
    For $\psi \in \End_{K}(V)$, we still use $\psi$ to denote
    $$1 \otimes \psi \in \End_{K[T]}(K[T]\otimes V).$$
    $$\begin{aligned}
        P_{g \circ \varphi \circ g^{-1}} = & \det \left(T \Id_{V} - g \circ \varphi \circ g_{-1}\right)\\
        = & \det\left(g \circ (T \Id_{V} - \varphi) \circ g^{-1}\right)\\
        = & \det(g) \cdot \det(T \Id_{V} - \varphi) \det(g^{-1})\\
        = & P_{\varphi} (T).
    \end{aligned}$$
    For any $F \in K[T]$,
    $$g \circ F(\varphi) \circ g^{-1} = F(g \circ \varphi \circ g^{-1})$$
    since 
    $$g \circ \varphi^n \circ g^{-1} = (g \circ \varphi \circ g^{-1})^n.$$
    Hence
    $$\begin{aligned}
        \{F \in K[T] \mid F(\varphi) = 0\} & = \{F \in K[T] \mid g \circ F(\varphi) \circ g^{-1} = 0\}\\
        & = \{F \in K[T] \mid F(g \circ \varphi \circ g^{-1}) = 0\}.
    \end{aligned}$$
    We obtain $Q_{\varphi} = Q_{g \circ \varphi \circ g^{-1}}$.
\end{proofenv}

\begin{theoremenv}
    Let $V$ and $W$ be finite dimensional vector spaces over $K$,
    $\varphi : V \rightarrow W$, $\psi : W \rightarrow V$ be $K$-linear mappings.
    Then
    $$P_{\psi \circ \varphi}(T) T^{\dim(W)} = P_{\varphi\circ\psi}(T) T^{\dim(V)}.$$
\end{theoremenv}
\begin{proofenv}
    Consider
    $$\begin{array}{rrcl}
        \Phi : & V \oplus W & \longrightarrow & V \oplus W\\
        & (x,y) & \longmapsto & (\psi(\varphi(x)), \varphi(x)).
    \end{array}$$
    Hence $P_{\Phi}(T) = P_{\psi \circ \varphi}(T) T^{\dim(W)} = P_{\psi\circ\varphi}(T) T^{\dim(V)}$.
    In fact, if $(x_i)_{i = 1}^n$ and $(y_{j})_{j = 1}^m$ are bases of $V$ and $W$, respectively, then
    $$\begin{aligned}
        & (T - \Phi)(x_1)\wedge \cdots \wedge (T-\Phi)(x_n) \wedge (T-\Phi)(y_1)\wedge \cdots \wedge (T-\Phi)(y_m)\\
        = & T^m(Tx_1 - \psi(\varphi(x_1))- \varphi(x_1)) \wedge \cdots\\
        & \wedge (Tx_n - \psi(\varphi(x_n)) - \varphi(x_n)) \wedge y_1 \wedge \cdots \wedge y_m\\
        = & T^m(Tx_1 - \psi(\varphi(x_1))) \wedge \cdots \wedge (Tx_n - \psi(\varphi(x_n))) \wedge y_1 \wedge \cdots \wedge y_m\\
        = & P_{\psi \circ \varphi}(T) T^{m} (x_1 \wedge \cdots \wedge x_n \wedge y_1 \wedge \cdots \wedge y_m).
    \end{aligned}.$$
    Let
    $$\begin{array}{rrcl}
        f : & V \oplus W & \longrightarrow & V \oplus W\\
        & (x,y) & \longmapsto & (x - \psi(y), y).
    \end{array}$$
    $f$ is invertible and, $f^{-1}(x,y) = (x + \psi(y),y)$,
    $$\begin{aligned}
        (f \circ \Phi \circ f^{-1})(x,y) = & (f \circ \Phi)(x + \psi(y), y)\\
        = & f(\psi(\varphi(x)) + \psi(\varphi(\psi(y))), \varphi(x) + \varphi(\psi(y)))\\
        = & (0, \varphi(x) + \varphi(\psi(y)))\\
    \end{aligned}$$
    $$\begin{aligned}
        & (T - f \circ \Phi \circ f^{-1})(x_1) \wedge \cdots \wedge (T - f \circ \Phi \circ f^{-1})(x_n)\\
        & \wedge (T - f \circ \Phi \circ f^{-1})(y_1) \wedge \cdots \wedge (T - f \circ \Phi \circ f^{-1})(y_m)\\
        = & (T x_1 - \varphi(x_1)) \wedge \cdots \wedge (T x_n - \varphi(x_n))\\
        & \wedge (T y_1 - \varphi(\psi(y_1))) \wedge \cdots \wedge (T y_m - \varphi(\psi(y_m)))\\
        = & P_{\varphi \circ \psi}(T) (T x_1 - \varphi(x_1)) \wedge \cdots \wedge (T x_n - \varphi(x_n)) \wedge y_1 \wedge \cdots \wedge y_m\\
        = & T^n P_{\varphi \circ \psi} (T).
    \end{aligned}$$
\end{proofenv}

\begin{corollaryenv}
    $$\tr(\varphi \circ \psi) = \tr(\psi \circ \varphi).$$
\end{corollaryenv}
\begin{proofenv}
    $$P_{\varphi \circ \psi}(T) = T^{\dim(W)} - \tr(\varphi \circ \psi) T^{\dim(W) - 1} + \cdots,$$
    $$P_{\psi \circ \varphi}(T) = T^{\dim(V)} - \tr(\psi \circ \varphi) T^{\dim(V) - 1} + \cdots.$$
\end{proofenv}
